\documentclass[12pt,a4paper]{article}
\usepackage[utf8]{vietnam}
\usepackage{amsmath,amssymb}
\usepackage{tabularx}
\usepackage{tikz}
\usepackage{xcolor}
\usepackage[top=0.6cm, bottom=0.6cm, left=0.6cm, right=0.6cm]{geometry}
\usepackage[dethi]{ex_test}

\SAOPTN{widthSA=2.8,heightSA=0.62}
\setlength{\parskip}{0pt}

% Dòng kẻ chấm mờ chuẩn viết tay
\newcommand{\dongke}[1]{%
  \begingroup%
  \setlength{\parskip}{0pt}%
  \baselineskip=18.5pt%
  \vspace{1.0mm}%
  \foreach \i in {1,...,#1}{%
    \noindent{\color{gray!60}\dotfill}\par%
  }%
  \vspace{0.6mm}%
  \endgroup%
}

% Định nghĩa ô bài tập tự luận trong bảng
\newcommand{\oBaiTap}[3][6]{%
  \vspace{1.0mm}%
  \noindent\textbf{#2} #3\par%
  \dongke{#1}%
}

\begin{document}
\pagestyle{empty}
% =============================================================================
% ========================= MA DE 101 - MAT TRUOC (TRANG 1/2) =============
% =============================================================================
\noindent
\begin{tabularx}{\linewidth}{|X|c|}
\hline
\textbf{PHIẾU KIỂM TRA ĐỊNH KỲ TOÁN 6 - TUẦN 10 - BUỔI 01} & \textbf{Mã đề: 101 - Trang 1/2} \\
\hline
\end{tabularx}

\vspace{1.0mm}

\noindent
\begin{tabularx}{\linewidth}{|p{3.2cm}|X|}
\hline
\centering\textbf{Điểm số} & \centering\textbf{Lời nhận xét của thầy cô} \tabularnewline
\rule{0pt}{1.8cm} & \\
\hline
\end{tabularx}

\vspace{1.0mm}

\noindent\textbf{Câu 1 (0,25 điểm).} Nhiệt độ $5^\circ\text{C}$ dưới $0^\circ\text{C}$ được biểu diễn bằng số nguyên là: \shortans{$-5^\circ\text{C}$}\par\vspace{0.25mm}
\noindent\textbf{Câu 2 (0,25 điểm).} Điền kí hiệu thích hợp $(\in, \notin)$ vào chỗ chấm: $-7 \;\dots\; \mathbb{N}$ \shortans{$\notin$}\par\vspace{0.25mm}
\noindent\textbf{Câu 3 (0,25 điểm).} Số đối của số $-18$ là: \shortans{$18$}\par\vspace{0.25mm}
\noindent\textbf{Câu 4 (0,25 điểm).} Điểm ở bên trái gốc $0$ và cách gốc $0$ một khoảng $4$ đơn vị là: \shortans{$-4$}\par\vspace{0.25mm}
\noindent\textbf{Câu 5 (0,25 điểm).} Khoảng cách giữa điểm $-3$ và điểm $5$ trên trục số bằng: \shortans{$8$}\par\vspace{0.25mm}
\noindent\textbf{Câu 6 (0,25 điểm).} Có bao nhiêu số nguyên $x$ thỏa mãn điều kiện $-3 \le x \le 4$? \shortans{$8$}\par\vspace{0.25mm}
\noindent\textbf{Câu 7 (0,25 điểm).} Tổng tất cả các số nguyên $x$ thỏa mãn $-5 < x < 5$ bằng: \shortans{$0$}\par\vspace{0.25mm}
\noindent\textbf{Câu 8 (0,25 điểm).} Hình có tâm đối xứng trong: Vuông, Tam giác đều, Thang cân là: \shortans{Hình vuông}\par\vspace{0.25mm}
\noindent\textbf{Câu 9 (0,25 điểm).} Hình bình hành có đáy $15\text{ cm}$, chiều cao tương ứng $8\text{ cm}$. Diện tích là: \shortans{$120\text{ cm}^2$}\par\vspace{0.25mm}
\noindent\textbf{Câu 10 (0,25 điểm).} Nhiệt độ Sa Pa sáng là $3^\circ\text{C}$, đêm giảm $5^\circ\text{C}$. Nhiệt độ ban đêm là: \shortans{$-2^\circ\text{C}$}\par\vspace{0.25mm}

\vspace{0.3mm}
\noindent\textbf{Câu 11 (3,0 điểm). Thực hiện phép tính (hợp lý nếu có thể):}
\vspace{0.4mm}

\noindent
\begin{tabularx}{\linewidth}{|X|X|}
\hline
\oBaiTap[7]{a)}{$36 + 148 + 64 + 52$} & \oBaiTap[7]{b)}{$46 \cdot 68 + 46 \cdot 32$} \\
\hline
\oBaiTap[7]{c)}{$7 \cdot 2^3 - 36 : 3^2 + 2026^0$} & \oBaiTap[7]{d)}{$160 : (28 - 12) + 4^2 \cdot 3$} \\
\hline
\oBaiTap[7]{e)}{$150 : [35 - (11 - 6)^2]$} & \oBaiTap[7]{f)}{$120 : \{50 - [35 - (10 - 5)^2]\}$} \\
\hline
\end{tabularx}
\newpage
% =============================================================================
% ========================= MA DE 101 - MAT SAU (TRANG 2/2) ==============
% =============================================================================
\noindent
\begin{tabularx}{\linewidth}{|X|}
\hline
\centering\textbf{TOÁN 6 - TUẦN 10 - BUỔI 01 | Mã đề: 101 - Trang 2/2} \\
\hline
\end{tabularx}

\vspace{1.2mm}

\noindent\textbf{Câu 12 (2,0 điểm). Tìm số tự nhiên / số nguyên $x$, biết:}
\vspace{0.5mm}

\noindent
\begin{tabularx}{\linewidth}{|X|X|}
\hline
\oBaiTap[5]{a)}{$3x - 17 = 28$} & \oBaiTap[5]{b)}{$5(x + 7) - 15 = 45$} \\
\hline
\oBaiTap[7]{c)}{Tập hợp các số nguyên $x$ sao cho $-3 \le x < 4$} & \oBaiTap[7]{d)}{$2 \cdot 3^{x-1} + 17 = 71$} \\
\hline
\end{tabularx}

\vspace{1.2mm}

\noindent
\begin{tabularx}{\linewidth}{|X|}
\hline
\vspace{0.5mm}
\textbf{Câu 13 (2,0 điểm).} Đỉnh núi Fansipan có độ cao $3\,143\text{ m}$ trên mực nước biển. Đáy Rãnh Mariana nằm ở độ sâu $10\,994\text{ m}$ dưới mực nước biển.\par
\textbf{a) (1,0 điểm)} Dùng số nguyên thích hợp để biểu diễn độ cao đỉnh Fansipan và độ sâu đáy Rãnh Mariana so với mực nước biển.\par
\textbf{b) (1,0 điểm)} Tính khoảng cách theo phương thẳng đứng (độ chênh lệch độ cao) giữa đỉnh Fansipan và đáy Rãnh Mariana.\par
\vspace{0.4mm}
\noindent\textit{Bài giải:}\dongke{11}
\\ \hline
\end{tabularx}

\vspace{1.2mm}

\noindent
\begin{tabularx}{\linewidth}{|X|}
\hline
\vspace{0.5mm}
\textbf{Câu 14 (0,5 điểm).} Tìm tất cả các số nguyên $n$ để phân số $A = \dfrac{n + 5}{n + 1}$ nhận giá trị là một số nguyên.\par
\vspace{0.4mm}
\noindent\textit{Bài giải:}\dongke{7}
\\ \hline
\end{tabularx}

\newpage
% =============================================================================
% ========================= MA DE 102 - MAT TRUOC (TRANG 1/2) =============
% =============================================================================
\noindent
\begin{tabularx}{\linewidth}{|X|c|}
\hline
\textbf{PHIẾU KIỂM TRA ĐỊNH KỲ TOÁN 6 - TUẦN 10 - BUỔI 01} & \textbf{Mã đề: 102 - Trang 1/2} \\
\hline
\end{tabularx}

\vspace{1.0mm}

\noindent
\begin{tabularx}{\linewidth}{|p{3.2cm}|X|}
\hline
\centering\textbf{Điểm số} & \centering\textbf{Lời nhận xét của thầy cô} \tabularnewline
\rule{0pt}{1.8cm} & \\
\hline
\end{tabularx}

\vspace{1.0mm}

\noindent\textbf{Câu 1 (0,25 điểm).} Độ sâu $20\text{ m}$ dưới mực nước biển được biểu diễn bằng số nguyên là: \shortans{$-20\text{ m}$}\par\vspace{0.25mm}
\noindent\textbf{Câu 2 (0,25 điểm).} Điền kí hiệu thích hợp $(\in, \notin)$ vào chỗ chấm: $-12 \;\dots\; \mathbb{Z}$ \shortans{$\in$}\par\vspace{0.25mm}
\noindent\textbf{Câu 3 (0,25 điểm).} Số đối của số $25$ là: \shortans{$-25$}\par\vspace{0.25mm}
\noindent\textbf{Câu 4 (0,25 điểm).} Điểm ở bên trái gốc $0$ và cách gốc $0$ một khoảng $6$ đơn vị là: \shortans{$-6$}\par\vspace{0.25mm}
\noindent\textbf{Câu 5 (0,25 điểm).} Khoảng cách giữa điểm $-4$ và điểm $6$ trên trục số bằng: \shortans{$10$}\par\vspace{0.25mm}
\noindent\textbf{Câu 6 (0,25 điểm).} Có bao nhiêu số nguyên $x$ thỏa mãn điều kiện $-4 \le x \le 5$? \shortans{$10$}\par\vspace{0.25mm}
\noindent\textbf{Câu 7 (0,25 điểm).} Tổng tất cả các số nguyên $x$ thỏa mãn $-6 < x < 6$ bằng: \shortans{$0$}\par\vspace{0.25mm}
\noindent\textbf{Câu 8 (0,25 điểm).} Chữ cái có cả trục và tâm đối xứng trong A, H, M, O là: \shortans{H, O}\par\vspace{0.25mm}
\noindent\textbf{Câu 9 (0,25 điểm).} Hình bình hành có đáy $16\text{ cm}$, chiều cao tương ứng $9\text{ cm}$. Diện tích là: \shortans{$144\text{ cm}^2$}\par\vspace{0.25mm}
\noindent\textbf{Câu 10 (0,25 điểm).} Tàu ngầm ở độ cao $-30\text{ m}$, lặn sâu thêm $20\text{ m}$. Độ cao mới của tàu là: \shortans{$-50\text{ m}$}\par\vspace{0.25mm}

\vspace{0.3mm}
\noindent\textbf{Câu 11 (3,0 điểm). Thực hiện phép tính (hợp lý nếu có thể):}
\vspace{0.4mm}

\noindent
\begin{tabularx}{\linewidth}{|X|X|}
\hline
\oBaiTap[7]{a)}{$47 + 136 + 53 + 64$} & \oBaiTap[7]{b)}{$57 \cdot 74 + 57 \cdot 26$} \\
\hline
\oBaiTap[7]{c)}{$8 \cdot 2^3 - 45 : 3^2 + 2025^0$} & \oBaiTap[7]{d)}{$180 : (35 - 17) + 5^2 \cdot 2$} \\
\hline
\oBaiTap[7]{e)}{$180 : [45 - (12 - 7)^2]$} & \oBaiTap[7]{f)}{$160 : \{50 - [35 - (12 - 7)^2]\}$} \\
\hline
\end{tabularx}
\newpage
% =============================================================================
% ========================= MA DE 102 - MAT SAU (TRANG 2/2) ==============
% =============================================================================
\noindent
\begin{tabularx}{\linewidth}{|X|}
\hline
\centering\textbf{TOÁN 6 - TUẦN 10 - BUỔI 01 | Mã đề: 102 - Trang 2/2} \\
\hline
\end{tabularx}

\vspace{1.2mm}

\noindent\textbf{Câu 12 (2,0 điểm). Tìm số tự nhiên / số nguyên $x$, biết:}
\vspace{0.5mm}

\noindent
\begin{tabularx}{\linewidth}{|X|X|}
\hline
\oBaiTap[5]{a)}{$4x - 19 = 37$} & \oBaiTap[5]{b)}{$4(x + 8) - 18 = 46$} \\
\hline
\oBaiTap[7]{c)}{Tập hợp các số nguyên $x$ sao cho $-4 < x \le 3$} & \oBaiTap[7]{d)}{$3 \cdot 2^{x+1} - 15 = 33$} \\
\hline
\end{tabularx}

\vspace{1.2mm}

\noindent
\begin{tabularx}{\linewidth}{|X|}
\hline
\vspace{0.5mm}
\textbf{Câu 13 (2,0 điểm).} Đỉnh núi Bà Đen có độ cao $986\text{ m}$ trên mực nước biển. Một tàu ngầm nghiên cứu đại dương đang hoạt động ở độ sâu $450\text{ m}$ dưới mực nước biển.\par
\textbf{a) (1,0 điểm)} Dùng số nguyên thích hợp để biểu diễn độ cao đỉnh núi Bà Đen và vị trí tàu ngầm so với mực nước biển.\par
\textbf{b) (1,0 điểm)} Tính khoảng cách theo phương thẳng đứng giữa đỉnh núi Bà Đen và vị trí tàu ngầm.\par
\vspace{0.4mm}
\noindent\textit{Bài giải:}\dongke{11}
\\ \hline
\end{tabularx}

\vspace{1.2mm}

\noindent
\begin{tabularx}{\linewidth}{|X|}
\hline
\vspace{0.5mm}
\textbf{Câu 14 (0,5 điểm).} Tìm tất cả các số nguyên $n$ để phân số $B = \dfrac{n + 7}{n + 2}$ nhận giá trị là một số nguyên.\par
\vspace{0.4mm}
\noindent\textit{Bài giải:}\dongke{7}
\\ \hline
\end{tabularx}

\newpage
% =============================================================================
% ========================= MA DE 103 - MAT TRUOC (TRANG 1/2) =============
% =============================================================================
\noindent
\begin{tabularx}{\linewidth}{|X|c|}
\hline
\textbf{PHIẾU KIỂM TRA ĐỊNH KỲ TOÁN 6 - TUẦN 10 - BUỔI 01} & \textbf{Mã đề: 103 - Trang 1/2} \\
\hline
\end{tabularx}

\vspace{1.0mm}

\noindent
\begin{tabularx}{\linewidth}{|p{3.2cm}|X|}
\hline
\centering\textbf{Điểm số} & \centering\textbf{Lời nhận xét của thầy cô} \tabularnewline
\rule{0pt}{1.8cm} & \\
\hline
\end{tabularx}

\vspace{1.0mm}

\noindent\textbf{Câu 1 (0,25 điểm).} Một doanh nghiệp bị lỗ $15$ triệu đồng được biểu diễn bằng số nguyên là: \shortans{$-15\text{ triệu}$}\par\vspace{0.25mm}
\noindent\textbf{Câu 2 (0,25 điểm).} Điền kí hiệu thích hợp $(\in, \notin)$ vào chỗ chấm: Số $0 \;\dots\; \mathbb{Z}^-$ (số nguyên âm) \shortans{$\notin$}\par\vspace{0.25mm}
\noindent\textbf{Câu 3 (0,25 điểm).} Số đối của số $-36$ là: \shortans{$36$}\par\vspace{0.25mm}
\noindent\textbf{Câu 4 (0,25 điểm).} Điểm ở bên trái gốc $0$ và cách gốc $0$ một khoảng $7$ đơn vị là: \shortans{$-7$}\par\vspace{0.25mm}
\noindent\textbf{Câu 5 (0,25 điểm).} Khoảng cách giữa điểm $-5$ và điểm $3$ trên trục số bằng: \shortans{$8$}\par\vspace{0.25mm}
\noindent\textbf{Câu 6 (0,25 điểm).} Có bao nhiêu số nguyên $x$ thỏa mãn điều kiện $-5 < x \le 3$? \shortans{$8$}\par\vspace{0.25mm}
\noindent\textbf{Câu 7 (0,25 điểm).} Tổng tất cả các số nguyên $x$ thỏa mãn $-4 < x \le 4$ bằng: \shortans{$4$}\par\vspace{0.25mm}
\noindent\textbf{Câu 8 (0,25 điểm).} Hình bình hành có tất cả bao nhiêu tâm đối xứng? \shortans{$1$}\par\vspace{0.25mm}
\noindent\textbf{Câu 9 (0,25 điểm).} Hình bình hành có đáy $18\text{ cm}$, chiều cao tương ứng $7\text{ cm}$. Diện tích là: \shortans{$126\text{ cm}^2$}\par\vspace{0.25mm}
\noindent\textbf{Câu 10 (0,25 điểm).} Tài khoản có $500\,000\text{ đ}$, rút $700\,000\text{ đ}$. Số dư tài khoản là: \shortans{$-200\,000\text{ đ}$}\par\vspace{0.25mm}

\vspace{0.3mm}
\noindent\textbf{Câu 11 (3,0 điểm). Thực hiện phép tính (hợp lý nếu có thể):}
\vspace{0.4mm}

\noindent
\begin{tabularx}{\linewidth}{|X|X|}
\hline
\oBaiTap[7]{a)}{$28 + 155 + 72 + 45$} & \oBaiTap[7]{b)}{$68 \cdot 63 + 68 \cdot 37$} \\
\hline
\oBaiTap[7]{c)}{$6 \cdot 2^3 - 27 : 3^2 + 2024^0$} & \oBaiTap[7]{d)}{$140 : (26 - 12) + 3^2 \cdot 6$} \\
\hline
\oBaiTap[7]{e)}{$210 : [55 - (13 - 8)^2]$} & \oBaiTap[7]{f)}{$200 : \{50 - [35 - (13 - 8)^2]\}$} \\
\hline
\end{tabularx}
\newpage
% =============================================================================
% ========================= MA DE 103 - MAT SAU (TRANG 2/2) ==============
% =============================================================================
\noindent
\begin{tabularx}{\linewidth}{|X|}
\hline
\centering\textbf{TOÁN 6 - TUẦN 10 - BUỔI 01 | Mã đề: 103 - Trang 2/2} \\
\hline
\end{tabularx}

\vspace{1.2mm}

\noindent\textbf{Câu 12 (2,0 điểm). Tìm số tự nhiên / số nguyên $x$, biết:}
\vspace{0.5mm}

\noindent
\begin{tabularx}{\linewidth}{|X|X|}
\hline
\oBaiTap[5]{a)}{$5x - 23 = 42$} & \oBaiTap[5]{b)}{$6(x + 5) - 24 = 48$} \\
\hline
\oBaiTap[7]{c)}{Tập hợp các số nguyên $x$ sao cho $-2 \le x \le 5$} & \oBaiTap[7]{d)}{$4 \cdot 3^{x-2} + 22 = 58$} \\
\hline
\end{tabularx}

\vspace{1.2mm}

\noindent
\begin{tabularx}{\linewidth}{|X|}
\hline
\vspace{0.5mm}
\textbf{Câu 13 (2,0 điểm).} Đỉnh núi Mẫu Sơn có độ cao $1\,541\text{ m}$ trên mực nước biển. Vực thẳm Puerto Rico có độ sâu $8\,376\text{ m}$ dưới mực nước biển.\par
\textbf{a) (1,0 điểm)} Dùng số nguyên thích hợp để biểu diễn độ cao đỉnh Mẫu Sơn và độ sâu đáy vực Puerto Rico so với mực nước biển.\par
\textbf{b) (1,0 điểm)} Tính khoảng cách theo phương thẳng đứng giữa đỉnh núi Mẫu Sơn và đáy vực Puerto Rico.\par
\vspace{0.4mm}
\noindent\textit{Bài giải:}\dongke{11}
\\ \hline
\end{tabularx}

\vspace{1.2mm}

\noindent
\begin{tabularx}{\linewidth}{|X|}
\hline
\vspace{0.5mm}
\textbf{Câu 14 (0,5 điểm).} Tìm tất cả các số nguyên $n$ để phân số $C = \dfrac{n + 8}{n + 2}$ nhận giá trị là một số nguyên.\par
\vspace{0.4mm}
\noindent\textit{Bài giải:}\dongke{7}
\\ \hline
\end{tabularx}

\newpage
% =============================================================================
% ========================= MA DE 104 - MAT TRUOC (TRANG 1/2) =============
% =============================================================================
\noindent
\begin{tabularx}{\linewidth}{|X|c|}
\hline
\textbf{PHIẾU KIỂM TRA ĐỊNH KỲ TOÁN 6 - TUẦN 10 - BUỔI 01} & \textbf{Mã đề: 104 - Trang 1/2} \\
\hline
\end{tabularx}

\vspace{1.0mm}

\noindent
\begin{tabularx}{\linewidth}{|p{3.2cm}|X|}
\hline
\centering\textbf{Điểm số} & \centering\textbf{Lời nhận xét của thầy cô} \tabularnewline
\rule{0pt}{1.8cm} & \\
\hline
\end{tabularx}

\vspace{1.0mm}

\noindent\textbf{Câu 1 (0,25 điểm).} Năm $287$ trước Công nguyên được biểu diễn bằng số nguyên là: \shortans{$-287$}\par\vspace{0.25mm}
\noindent\textbf{Câu 2 (0,25 điểm).} Điền kí hiệu thích hợp $(\subset, \not\subset)$ vào chỗ chấm: Tập hợp $\mathbb{N} \;\dots\; \mathbb{Z}$ \shortans{$\subset$}\par\vspace{0.25mm}
\noindent\textbf{Câu 3 (0,25 điểm).} Số đối của số $0$ là: \shortans{$0$}\par\vspace{0.25mm}
\noindent\textbf{Câu 4 (0,25 điểm).} Điểm ở bên trái gốc $0$ và cách gốc $0$ một khoảng $9$ đơn vị là: \shortans{$-9$}\par\vspace{0.25mm}
\noindent\textbf{Câu 5 (0,25 điểm).} Khoảng cách giữa điểm $-6$ và điểm $4$ trên trục số bằng: \shortans{$10$}\par\vspace{0.25mm}
\noindent\textbf{Câu 6 (0,25 điểm).} Có bao nhiêu số nguyên $x$ thỏa mãn điều kiện $-2 \le x < 7$? \shortans{$9$}\par\vspace{0.25mm}
\noindent\textbf{Câu 7 (0,25 điểm).} Tổng tất cả các số nguyên $x$ thỏa mãn $-5 \le x < 5$ bằng: \shortans{$-5$}\par\vspace{0.25mm}
\noindent\textbf{Câu 8 (0,25 điểm).} Hình chữ nhật có tất cả bao nhiêu trục đối xứng? \shortans{$2$}\par\vspace{0.25mm}
\noindent\textbf{Câu 9 (0,25 điểm).} Hình bình hành có đáy $20\text{ cm}$, chiều cao tương ứng $6\text{ cm}$. Diện tích là: \shortans{$120\text{ cm}^2$}\par\vspace{0.25mm}
\noindent\textbf{Câu 10 (0,25 điểm).} Kho đông lạnh ở $-8^\circ\text{C}$, nhân viên giảm thêm $4^\circ\text{C}$. Nhiệt độ mới là: \shortans{$-12^\circ\text{C}$}\par\vspace{0.25mm}

\vspace{0.3mm}
\noindent\textbf{Câu 11 (3,0 điểm). Thực hiện phép tính (hợp lý nếu có thể):}
\vspace{0.4mm}

\noindent
\begin{tabularx}{\linewidth}{|X|X|}
\hline
\oBaiTap[7]{a)}{$39 + 143 + 61 + 57$} & \oBaiTap[7]{b)}{$79 \cdot 81 + 79 \cdot 19$} \\
\hline
\oBaiTap[7]{c)}{$9 \cdot 2^3 - 54 : 3^2 + 2027^0$} & \oBaiTap[7]{d)}{$200 : (38 - 18) + 4^2 \cdot 4$} \\
\hline
\oBaiTap[7]{e)}{$240 : [65 - (14 - 9)^2]$} & \oBaiTap[7]{f)}{$240 : \{50 - [35 - (14 - 9)^2]\}$} \\
\hline
\end{tabularx}
\newpage
% =============================================================================
% ========================= MA DE 104 - MAT SAU (TRANG 2/2) ==============
% =============================================================================
\noindent
\begin{tabularx}{\linewidth}{|X|}
\hline
\centering\textbf{TOÁN 6 - TUẦN 10 - BUỔI 01 | Mã đề: 104 - Trang 2/2} \\
\hline
\end{tabularx}

\vspace{1.2mm}

\noindent\textbf{Câu 12 (2,0 điểm). Tìm số tự nhiên / số nguyên $x$, biết:}
\vspace{0.5mm}

\noindent
\begin{tabularx}{\linewidth}{|X|X|}
\hline
\oBaiTap[5]{a)}{$6x - 27 = 45$} & \oBaiTap[5]{b)}{$3(x + 9) - 16 = 47$} \\
\hline
\oBaiTap[7]{c)}{Tập hợp các số nguyên $x$ sao cho $-5 < x < 3$} & \oBaiTap[7]{d)}{$5 \cdot 2^{x-1} - 13 = 27$} \\
\hline
\end{tabularx}

\vspace{1.2mm}

\noindent
\begin{tabularx}{\linewidth}{|X|}
\hline
\vspace{0.5mm}
\textbf{Câu 13 (2,0 điểm).} Đỉnh núi Pu Si Lung cao $3\,083\text{ m}$ trên mực nước biển. Bề mặt Biển Chết nằm ở độ sâu $430\text{ m}$ dưới mực nước biển.\par
\textbf{a) (1,0 điểm)} Dùng số nguyên thích hợp để biểu diễn độ cao đỉnh Pu Si Lung và bề mặt Biển Chết so với mực nước biển.\par
\textbf{b) (1,0 điểm)} Tính khoảng cách theo phương thẳng đứng giữa đỉnh Pu Si Lung và bề mặt Biển Chết.\par
\vspace{0.4mm}
\noindent\textit{Bài giải:}\dongke{11}
\\ \hline
\end{tabularx}

\vspace{1.2mm}

\noindent
\begin{tabularx}{\linewidth}{|X|}
\hline
\vspace{0.5mm}
\textbf{Câu 14 (0,5 điểm).} Tìm tất cả các số nguyên $n$ để phân số $D = \dfrac{n + 6}{n + 3}$ nhận giá trị là một số nguyên.\par
\vspace{0.4mm}
\noindent\textit{Bài giải:}\dongke{7}
\\ \hline
\end{tabularx}
\end{document}
